\documentclass[11pt,a4paper]{report}
\usepackage[a4paper,margin=1in]{geometry}
\usepackage[T1]{fontenc}
\usepackage{lmodern}
\usepackage{microtype}
\usepackage{booktabs}
\usepackage{amsmath,amssymb}
\usepackage{graphicx}
\usepackage{hyperref}
\usepackage{xcolor}
\usepackage{enumitem}
\usepackage{caption}
\hypersetup{colorlinks=true, linkcolor=black, citecolor=black, urlcolor=blue}

\title{Privacy-Preserving Federated Learning for IoT-Based Healthcare:\\
A Layered-Defense Benchmark on ECG Arrhythmia Classification}
\author{Mapeke}
\date{2026}

\begin{document}

\maketitle

\begin{abstract}
\noindent
Federated learning (FL) is the textbook answer to privacy-preserving training on
healthcare IoT data, and the layered defenses around it --- DP-SGD, secure
aggregation, and update compression --- are textbook ingredients of a
``privacy-preserving FL'' pipeline. The textbook does not say which layer pays
for which threat, or what each one costs in detection utility on a realistic
clinical workload. This thesis builds a fully working layered-defense framework
in roughly 2{,}500 lines of Python (Flower~+~PyTorch~+~Opacus), and benchmarks it
on the MIT-BIH Arrhythmia dataset with four simulated wearable IoT clients
training a 1-D CNN over the AAMI 5-class label scheme. Three findings drive the
thesis: (1) federation alone is essentially free --- vanilla FedAvg matches a
centralized baseline within 0.02\,pp test accuracy on this workload; (2) pairwise-mask
secure aggregation is also free --- it shifts test accuracy by 0.20\,pp, well below
seed variance, and the cryptographic correctness property (mask cancellation under
summation) is encoded as a unit test that runs in CI; (3) DP-SGD at $\varepsilon = 5$
collapses the model to the majority-class predictor on this heavily imbalanced
workload, with final accuracy equal to the prevalence of class \emph{N} to
two decimals --- a privacy/utility cliff that is invisible if one only
reports overall accuracy. Adding compression on top of DP does not degrade
the model further, because there is nothing left to degrade.
A follow-up run on MIT-BIH records 200--234 (where the imbalance is moderate,
82\,\% N rather than 97\,\%) shows that this collapse is dataset-conditional:
DP-SGD alone retains a $+6.37$\,pp signal over the majority baseline on that
slice, but the full pipeline (DP + SecAgg + compression at $\varepsilon = 8$)
still collapses to majority. The cliff is real; its location is a property of
the workload's class balance.
Every number in the thesis is reproducible from a committed YAML configuration via
\texttt{python -m experiments.run\_*}.
\end{abstract}

\tableofcontents

\chapter{Introduction}\label{ch:intro}

\section{Motivation}

Connected health devices --- smartwatches, Holter monitors, bedside telemetry, the
implantable devices that already exist in patients --- produce streams of
physiological data that are simultaneously \emph{clinically valuable} and
\emph{highly sensitive}. Existing deployments centralize that data, which (i)
creates a single point of failure for breaches, (ii) exposes patients to
re-identification even after pseudonymization, and (iii) conflicts with the
data-minimization principle of HIPAA and GDPR. Federated learning was proposed
\cite{mcmahan2017fedavg} precisely to break this pattern: each device trains
locally on its own raw data, the central server only ever sees model updates.

The naive version of this story is incomplete in two ways that the deep-learning
literature is comfortable papering over:

\begin{enumerate}
  \item \textbf{Model updates leak.} Per-client gradients reconstruct training
  records under known attacks; a coalition of clients or an honest-but-curious
  server can recover patient data from the updates alone. Differential privacy
  \cite{abadi2016dpsgd} and secure aggregation \cite{bonawitz2017secagg} are the
  two canonical defenses, but each addresses a different threat and each
  introduces its own utility tax.
  \item \textbf{Healthcare data is heavily imbalanced.} ECG arrhythmia
  classification on MIT-BIH is roughly 97\,\% normal beats; rare classes can be
  $<\!0.1\,\%$. Privacy mechanisms designed and tuned on balanced datasets
  (CIFAR, MNIST) behave badly here in ways the headline-accuracy literature
  rarely surfaces.
\end{enumerate}

This thesis takes both of these seriously. We build a working layered-defense
framework, run it end-to-end on real PhysioNet data, and report what each
privacy layer actually costs --- including the cases where the cost is the
detector ceasing to detect anything.

\section{Contributions}

\textbf{C1 --- A layered-defense FL framework with separable privacy layers.}
The Flower client interface used here exposes three independent switches
(differential privacy, secure aggregation, compression) which can be enabled in
isolation or jointly via a single YAML config. The marginal cost of each layer
on test accuracy, spent $\varepsilon$, and convergence is measured by holding
the others fixed (Chapters~\ref{ch:detection},~\ref{ch:secagg},~\ref{ch:dp}).
The implementation is roughly 2{,}500 lines of Python and runs end-to-end on a
single workstation in under 30 minutes.

\textbf{C2 --- Cryptographic correctness as a code-level invariant.}
The pairwise-mask secure aggregation protocol's correctness argument
(every mask appears with $+$ and $-$ exactly once, so summing client updates
cancels every mask) is encoded as a unit test
(\texttt{tests/test\_secagg.py::test\_masks\_cancel\_under\_aggregation}) that
runs in CI. The argument from
\texttt{docs/privacy-analysis.md~\S2}
\cite{bonawitz2017secagg} is therefore mechanically checked rather than
prose-only.

\textbf{C3 --- A class-prior collapse diagnostic for DP on imbalanced
healthcare data.} On the MIT-BIH workload used here, the DP-SGD run at
$\varepsilon \approx 5$ converges to a final test accuracy of
\textbf{0.9711}, which equals the prevalence of class \emph{N} (Normal) to four
decimals. The full-pipeline run (DP + SecAgg + compression) at
$\varepsilon \approx 8$ converges to the \emph{same} 0.9711. We argue
(Chapter~\ref{ch:dp}, \S\ref{sec:f6.1}) that this is a class-prior collapse,
not a learning curve, and that it is the single most consequential and least
discussed failure mode of DP-SGD on realistic clinical data. A follow-up on
records 200--234 (\S\ref{sec:f6.3}) confirms the diagnostic generalises but
its severity is class-balance-conditional: on moderate imbalance DP-SGD
retains a measurable rare-class signal, while the full pipeline still
collapses.

\section{Structure}

Chapter~\ref{ch:bg} reviews FL, the three privacy primitives we layer, and the
MIT-BIH dataset.
Chapter~\ref{ch:related} situates the framework in the existing literature on
FL aggregation, DP-SGD, SecAgg, gradient compression, privacy-preserving FL
on healthcare data, and the membership-inference and gradient-inversion
attacks that motivate each layer.
Chapter~\ref{ch:threat} formalises the threat model: the assets being
protected, five adversaries A1--A5 with their capabilities and goals, the
attack-surface coverage matrix, and the trust assumptions the framework
inherits.
Chapter~\ref{ch:method} describes the implementation:
\texttt{FlowerHealthcareClient}, \texttt{PrivacyAwareFedAvg},
\texttt{DPConfig+attach\_dp}, the pairwise-mask SecAgg primitives, and the
top-$k$/int8 compression pipeline. Chapter~\ref{ch:detection} reports the
centralized baseline and vanilla FedAvg. Chapter~\ref{ch:secagg} adds secure
aggregation. Chapter~\ref{ch:dp} adds differential privacy and the full
pipeline. Chapter~\ref{ch:disc} discusses anticipated reviewer objections
and limitations. Chapter~\ref{ch:concl} concludes.

\section{Scope and non-goals}

\begin{itemize}[leftmargin=*]
  \item \textbf{We do not propose a new privacy primitive.} The argument is
  empirical: existing primitives composed under a protocol that makes their
  individual costs visible.
  \item \textbf{We do not deploy on real IoT hardware.} Every number is from a
  Flower simulation on a single GPU. The IoT angle is in the data partitioning
  (per-patient non-IID) and in the device-profile metadata (\texttt{wearable},
  \texttt{bedside}, \texttt{gateway}); the experiments do not exercise actual
  network or compute constraints.
  \item \textbf{We do not evaluate against poisoning / Byzantine clients.} The
  threat model is honest-but-curious server plus an external adversary holding
  the released global model (Appendix \texttt{docs/threat-model.md}). Robust
  aggregation (Krum, median, trimmed mean) is listed as future work
  (\S\ref{sec:future}).
  \item \textbf{Our primary experiments use MIT-BIH records 100--115},
  which \texttt{wfdb.dl\_database} pulls down by default and which exhibit
  the heavy-imbalance regime (97.1\,\% N) that drives the central finding.
  A follow-up section (Chapter~\ref{ch:dp}, \S\ref{sec:f6.3}) re-runs all
  five experiments on records 200--234 (more arrhythmia variety, 82.1\,\% N)
  to test whether the DP class-prior collapse is dataset-conditional.
\end{itemize}

\chapter{Background}\label{ch:bg}

\section{Federated learning and FedAvg}

The federated learning protocol of McMahan et~al.~\cite{mcmahan2017fedavg}
distributes training across $N$ clients each holding a local dataset $D_i$,
with the server orchestrating $T$ rounds of:
$$
w_{t+1} \;=\; \sum_{i=1}^{N} \tfrac{n_i}{n}\, \mathrm{LocalSGD}(w_t, D_i),
$$
where $n_i = |D_i|$ and $n = \sum_i n_i$. Raw $D_i$ never leaves client $i$;
only $\mathrm{LocalSGD}(w_t, D_i)$ does. FedAvg is the baseline we layer
defenses on top of.

The healthcare specialization is straightforward: each $D_i$ is a per-patient
or per-device partition of physiological data. We use four simulated wearable
clients holding disjoint Dirichlet-non-IID partitions
(\S\ref{sec:method-data}) of MIT-BIH ECG beats.

\section{The privacy stack}

Three layers, each addressing a different threat.

\paragraph{Differential privacy (DP-SGD).}
Replaces every local SGD step with: clip per-sample gradients to
$\ell_2$-norm $C$, add Gaussian noise of scale $\sigma C$, then step.
Composing across $T$ rounds and converting to $(\varepsilon, \delta)$ via the
RDP accountant \cite{mironov2017rdp} gives a tight bound on what an
adversary holding only the released global model can learn about any
single training record. Implemented via Opacus \cite{yousefpour2021opacus}.

\paragraph{Pairwise-mask secure aggregation.}
For every unordered client pair $\{i,j\}$, both clients derive a shared mask
vector $m_{ij}$. Client $i$ adds $+m_{ij}$ to its update; client $j$ adds
$-m_{ij}$. The server's sum cancels every mask exactly, so the server only ever
sees the aggregate $\sum_i \Delta w_i$ and never any individual $\Delta w_i$
\cite{bonawitz2017secagg}. Defends against an honest-but-curious server that
would otherwise inspect per-client updates.

\paragraph{Top-$k$ + int8 compression.}
Each client keeps the largest-magnitude $k$ fraction of its update entries
and quantizes them into 8-bit integers. Reduces bandwidth roughly $40\times$
for $k = 0.1$. Post-processing-immune to DP and compatible with SecAgg as
long as both endpoints agree on the sparsification mask.

\paragraph{Composition.}
DP and SecAgg compose without interaction: the DP guarantee is on $w_T$
regardless of channel; SecAgg is on intermediate per-round messages.
Compression composes with both as a post-hoc lossy operator. The full
composition argument is laid out in \S\ref{sec:composition}.

\medskip
The next three subsections give the formal statement of each layer:
\S\ref{sec:dpsgd-formal} writes out the DP-SGD step and the RDP-to-$(\varepsilon,
\delta)$ conversion; \S\ref{sec:secagg-formal} writes out the
mask-cancellation algebra of pairwise-mask SecAgg; \S\ref{sec:composition}
states how the three layers compose. The framework's design choices in
Chapter~\ref{ch:method} are direct implementations of these statements, and
the reproducibility check in Chapter~\ref{ch:detection} is built on top of
them.

\subsection{DP-SGD formal definition}\label{sec:dpsgd-formal}

Following Abadi et~al.~\cite{abadi2016dpsgd}, every local SGD step on a
minibatch $B$ of size $b$ is replaced by:

\begin{enumerate}[leftmargin=*]
  \item Compute per-sample gradients
        $g_i = \nabla_w \ell(w; x_i, y_i)$ for each $(x_i, y_i) \in B$.
  \item Clip per-sample gradients to $\ell_2$-norm at most $C$:
        $\bar g_i = g_i \,/\, \max\!\left(1, \tfrac{\|g_i\|_2}{C}\right)$.
        This bounds the sensitivity of the per-sample sum to any single
        training record by $C$.
  \item Add Gaussian noise of scale $\sigma C$ to the clipped sum and
        average:
        $\tilde g \;=\; \tfrac{1}{b}\!\left(\sum_{i \in B} \bar g_i \;+\;
         \mathcal{N}(0,\, \sigma^2 C^2 I)\right)$.
  \item Update weights: $w \leftarrow w - \eta \,\tilde g$.
\end{enumerate}

Each step satisfies $(\alpha,\, \alpha C^2 \sigma^{-2}/2)$-Rényi differential
privacy at every order $\alpha > 1$~\cite{mironov2017rdp}. Composing $T$
steps under RDP and converting to $(\varepsilon, \delta)$-DP via
$\varepsilon = \alpha C^2 \sigma^{-2} T / 2 + \log(1/\delta)/(\alpha - 1)$
(optimised over $\alpha$) gives the tight cumulative bound that Opacus's
\texttt{RDPAccountant}~\cite{yousefpour2021opacus} reports as
\texttt{get\_epsilon(steps)}. In the framework the user supplies
$(\varepsilon_{\mathrm{target}}, \delta_{\mathrm{target}}, C, T)$ and Opacus
auto-calibrates $\sigma$ to hit $\varepsilon \le \varepsilon_{\mathrm{target}}$
over $T$ steps; the framework reports the realised $\varepsilon$ at the end
of training (see Chapter~\ref{ch:dp}).

\paragraph{Parameter choices.} Table~\ref{tab:dp-defaults} lists the
defaults used throughout the thesis. \texttt{target\_epsilon} = 5 is a
conventional ``moderate'' budget for medical models in the literature;
\texttt{target\_delta} = $10^{-5}$ satisfies $\delta < 1/N$ for our dataset
size; $C = 1.0$ is the Opacus default and is well-matched to the gradient
norms we observe during training; $\sigma$ is auto-calibrated.

\begin{table}[h]
  \centering
  \caption{DP-SGD defaults used in the thesis. The user-facing knobs are
  $(\varepsilon_{\mathrm{target}}, \delta_{\mathrm{target}}, C)$; the
  noise multiplier $\sigma$ is solved for by the RDP accountant.}
  \label{tab:dp-defaults}
  \begin{tabular}{lll}
    \toprule
    Parameter & Default & Justification \\
    \midrule
    \texttt{target\_epsilon}     & 5.0    & Conventional moderate budget for medical models \\
    \texttt{target\_delta}       & $10^{-5}$ & $< 1/N$ for $N \approx 25{,}000$ training records \\
    \texttt{max\_grad\_norm} ($C$) & 1.0  & Opacus default; matches observed gradient scale \\
    \texttt{noise\_multiplier} ($\sigma$) & auto & Solved by RDP accountant to hit \texttt{target\_epsilon} \\
    \texttt{accountant}          & rdp    & RDP gives the tightest known bounds for Gaussian noise \\
    \bottomrule
  \end{tabular}
\end{table}

\paragraph{What $\varepsilon$ buys, intuitively.}
For the worst-case adversary A4 (Chapter~\ref{ch:threat}, \S3.2.4) with
auxiliary information about every other training record, the posterior on
whether a given record was in the training set shifts by at most a factor
of $e^\varepsilon$: at $\varepsilon = 5$ this is roughly 148, a loose
worst-case bound. Empirical membership-inference accuracy at this budget
is typically much closer to the 50\% random-guess baseline. The thesis
does not run a membership-inference attack to verify this empirically
because the $(\varepsilon, \delta)$-DP guarantee is a worst-case
\emph{distributional} statement that does not require empirical
attestation; running the attack would, at best, lower-bound something the
budget already upper-bounds.

\subsection{Pairwise-mask SecAgg correctness}\label{sec:secagg-formal}

Let each client $i \in [N]$ hold a parameter update $\Delta w_i \in
\mathbb{R}^d$ at the end of a training round. For every unordered pair
$\{i, j\}$ the two clients derive a shared seed $s_{ij}$ from a
key-exchange protocol (Diffie--Hellman in deployment-grade
implementations~\cite{bonawitz2017secagg}; a pre-shared secret in the
prototype, see Chapter~\ref{ch:method}, \S\ref{sec:method-strategy}). Each
client expands $s_{ij}$ via a cryptographic PRG into a mask vector
$m_{ij} \in \mathbb{R}^d$. Client $i$'s outgoing message is then:
\begin{equation}\label{eq:masked-update}
  \tilde{\Delta w}_i \;=\; \Delta w_i \;+\; \sum_{j > i} m_{ij}
                                   \;-\; \sum_{j < i} m_{ji} \pmod{p},
\end{equation}
where the modular reduction is over a fixed-point representation of
$\mathbb{R}^d$ in finite-field arithmetic (modulus $p$ in deployment-grade
implementations; \texttt{float32} arithmetic in the prototype, an accepted
simplification with documented numerical consequences).

\paragraph{Cancellation.} Summing
Equation~\eqref{eq:masked-update} across all clients:
\begin{align}
  S \;&=\; \sum_{i=1}^{N} \tilde{\Delta w}_i \nonumber \\
       &=\; \sum_{i=1}^{N} \Delta w_i
            \;+\; \sum_{i=1}^{N} \!\left( \sum_{j > i} m_{ij}
                                        - \sum_{j < i} m_{ji} \right) \nonumber \\
       &=\; \sum_{i=1}^{N} \Delta w_i
            \;+\; \sum_{\{i,j\}: i<j} (m_{ij} - m_{ij}) \nonumber \\
       &=\; \sum_{i=1}^{N} \Delta w_i. \label{eq:cancellation}
\end{align}
The second-to-last equality uses the fact that for every unordered pair
$\{i,j\}$ the mask $m_{ij}$ appears once with $+$ in client $i$'s
message (the term in $\sum_{j > i}$) and once with $-$ in client $j$'s
message (the term in $\sum_{j' < j}$ with $j' = i$). Every mask therefore
cancels exactly once at the sum level, regardless of $N$.

\paragraph{Privacy.} For any \emph{strict} subset $S \subsetneq [N]$ of
clients the marginal distribution of $\{\tilde{\Delta w}_i\}_{i \in S}$ is
uniformly random conditional on the true updates: at least one mask
$m_{ij}$ in the sum has $j \notin S$, that mask is uniformly random and
known only to the two endpoints, and adding a uniformly random vector to
any $\Delta w_i$ produces a uniformly random vector. The server therefore
learns no information about any individual $\Delta w_i$ beyond what is
already revealed by the total sum $\sum_i \Delta w_i$ that it must learn
to perform aggregation. The formal statement is Bonawitz et~al.\ Theorem
4.1~\cite{bonawitz2017secagg}.

\paragraph{Limitations of the prototype.}
The prototype departs from Bonawitz et~al.\ in three ways, each accepted
as a simulation-grade simplification and flagged in the source:
\textbf{(i)} pre-shared seeds per pair instead of Diffie--Hellman key
exchange — convenient for the simulation orchestrator, infeasible for a
real deployment because the seeds must be transported some other way;
\textbf{(ii)} no Shamir secret-sharing of seeds — the protocol does not
tolerate dropouts, since a missing client leaves un-cancelled masks in
the sum; \textbf{(iii)} \texttt{float32} arithmetic instead of modular
integer arithmetic over a finite field — introduces a small numerical
residue in the cancelled sum (\S\ref{sec:f5.2} measures this empirically
at 0.20\,pp end-to-end accuracy gap). All three are listed as
deployment-grade upgrades in \S\ref{sec:future}.

\subsection{Composition: DP + SecAgg + compression}\label{sec:composition}

The three layers compose without negotiation, in the precise sense that
the privacy guarantee of each layer survives the addition of the others:

\paragraph{DP guarantee survives SecAgg.} Differential privacy is a
guarantee on the released global model $w_T$ as a function of the input
training set. SecAgg changes the channel through which intermediate
updates flow, but does not change $w_T$ itself: the cancellation property
(Equation~\ref{eq:cancellation}) means the server reconstructs the true
$\sum_i \Delta w_i$ and applies it exactly as it would have without
SecAgg. The DP-SGD calibration of $\sigma$ to $\varepsilon_{\mathrm{target}}$
is therefore unaffected.

\paragraph{SecAgg guarantee survives DP.} SecAgg's privacy property is on
the per-round messages $\tilde{\Delta w}_i$: their joint distribution
conditional on the true updates is uniformly random for any strict
subset. DP changes the distribution of the true updates $\Delta w_i$ (by
adding Gaussian noise during local training), but does not change the
masking layer applied on top. SecAgg's guarantee composes with whatever
local update generation procedure the client runs.

\paragraph{Compression as post-processing.} Top-$k$ sparsification and
int8 quantisation are deterministic, lossy operators applied to
$\Delta w_i$ \emph{before} SecAgg masking. They do not weaken the DP
guarantee on $w_T$ by post-processing immunity (a fundamental property
of DP: any function of a DP output is itself DP at the same budget).
They do not weaken the SecAgg property either, because the SecAgg layer
operates on the post-compression vector and the cancellation algebra is
agnostic to the values being masked. The only requirement is that
sender and receiver agree on the sparsification mask; the framework
sends the indices alongside the values
(Chapter~\ref{ch:method}, \S\ref{sec:method-strategy}).

\paragraph{Composed statement.} Putting the three layers together:

\begin{quote}\itshape
  An adversary observing only the released global model $w_T$ learns at
  most $(\varepsilon, \delta)$ worth of information about any individual
  training record (\textbf{DP}), AND an adversary observing the per-round
  server inbox learns nothing about any individual client's update beyond
  the aggregate $\sum_i \Delta w_i$ (\textbf{SecAgg}). Compression reduces
  the per-round bandwidth from $4d \cdot 32\,\mathrm{bits}$ to roughly
  $4(k d \cdot 8\,\mathrm{bits} + k d \cdot 32\,\mathrm{bits})$ for
  top-$k$ + int8 at sparsity $k = 0.1$, a $\sim$40$\times$ reduction,
  without weakening either of the above guarantees.
\end{quote}

This is the contract that Chapters~\ref{ch:detection}--\ref{ch:dp}
benchmark empirically.

\section{Dataset and use case}

\paragraph{MIT-BIH Arrhythmia.}
PhysioNet's MIT-BIH database \cite{moody2001mitbih} is the de facto benchmark
for ECG arrhythmia classification. Sampling rate 360\,Hz, two leads,
beat-level annotations on $\sim$48 records. We use the first 12 records (100,
101, 103, 105, 106, 108, 109, 111, 112, 113, 114, 115), which
\texttt{wfdb.dl\_database} pulls without authentication.

\paragraph{AAMI 5-class scheme.}
Annotation symbols are mapped to the AAMI \cite{aami1998ec57} 5-class
grouping: \emph{N} (normal, including L/R/e/j), \emph{S} (supraventricular
ectopic), \emph{V} (ventricular ectopic), \emph{F} (fusion), \emph{Q}
(unknown / paced).

\paragraph{Class imbalance.}
Table~\ref{tab:class-counts} gives the realised distribution. The dataset
is dominated by class \emph{N}: it accounts for 97.1\,\% of the beats. Two
classes (F, Q) have under 10 examples each; no clinical claim about them is
possible from this data, but they are kept in the label space because they
appear in the AAMI scheme.

\begin{table}[h]
  \centering
  \caption{AAMI class distribution across the 12 MIT-BIH records, after
  beat-window segmentation (window size 360 samples). Class \emph{N}
  dominates with 97.1\,\% prevalence. Imbalance of this magnitude is the
  central technical fact of Chapter~\ref{ch:dp}.}
  \label{tab:class-counts}
  \begin{tabular}{lrr}
    \toprule
    Class & Count & Share \\
    \midrule
    \emph{N} (normal)               & 24{,}655 & 97.10\,\% \\
    \emph{V} (ventricular ectopic)  &     661 &  2.60\,\% \\
    \emph{S} (supraventricular)     &      62 &  0.24\,\% \\
    \emph{F} (fusion)               &       8 &  0.03\,\% \\
    \emph{Q} (unknown / paced)      &       7 &  0.03\,\% \\
    \midrule
    Total                           & 25{,}393 & 100.00\,\% \\
    \bottomrule
  \end{tabular}
\end{table}

\paragraph{Why ECG, why MIT-BIH.}
ECG monitoring is the canonical wearable-IoT health workload; MIT-BIH is the
canonical benchmark for it; the AAMI mapping is the canonical label scheme.
A privacy-preserving FL framework that ignores the imbalance characteristic of
this canonical workload is over-claiming.

\chapter{Related work}\label{ch:related}

This thesis contributes a layered-defence FL framework rather than a new
privacy primitive (Chapter~\ref{ch:intro}, \S1.4); the contribution rests on
how existing primitives compose and what each layer costs on a realistic
healthcare workload. This chapter places each of those primitives, and the
attacks they are built to mitigate, in the context of the existing
literature. Six strands matter: the FL aggregation rule that everything
sits on top of (\S\ref{sec:rw-fl}); the DP-SGD lineage that drives the
privacy guarantee on the released model (\S\ref{sec:rw-dp}); the secure
aggregation lineage that closes the per-update channel
(\S\ref{sec:rw-secagg}); the gradient compression literature that supplies
the bandwidth optimisation (\S\ref{sec:rw-compress}); prior
privacy-preserving FL work specifically on healthcare data
(\S\ref{sec:rw-health}); and the gradient-inversion and
membership-inference attacks that justify the threat model
(\S\ref{sec:rw-attacks}).

\section{Federated learning aggregation}\label{sec:rw-fl}

\paragraph{FedAvg.} The federated averaging protocol of McMahan
et~al.~\cite{mcmahan2017fedavg} introduced the now-canonical pattern: each
client runs $E$ local epochs of SGD on its private data, the server
weighted-averages the resulting parameter updates, and the protocol
iterates for $T$ communication rounds. FedAvg works well when client data
distributions are similar and the model is not over-parameterised. It is
the baseline of this thesis; \S\ref{sec:f4.1} confirms that on the MIT-BIH
workload it matches the centralized baseline within seed variance, so no
stronger aggregation rule is needed for the framework's contribution to be
visible.

\paragraph{FedProx.} Li et~al.~\cite{li2020fedprox} added a proximal
regulariser $\tfrac{\mu}{2}\|w_i - w_t\|^2$ to each client's local loss to
mitigate client drift in heterogeneous settings (different local data
distributions, different local-epoch counts). FedProx is strictly more
robust than FedAvg when the heterogeneity is severe but pays a small
constant cost otherwise. We do not use it because FedAvg already achieves
the centralized baseline.

\paragraph{SCAFFOLD.} Karimireddy et~al.\ proposed SCAFFOLD, which
estimates and corrects for client drift via per-client control variates,
producing strictly better convergence than FedAvg in non-IID settings. The
correction adds bookkeeping overhead and additional state per client.
Like FedProx, it is overkill for our workload.

\paragraph{Surveys.} Kairouz et~al.~\cite{kairouz2021advances} provide the
canonical survey of FL open problems and design space, including
privacy-preserving variants. We treat their taxonomy as the reference for
deployment-grade FL terminology.

\section{Differential privacy in deep learning}\label{sec:rw-dp}

\paragraph{Foundations.} Differential privacy was introduced by
Dwork~\cite{dwork2006dp} as a quantitative bound on what an adversary
holding the output of an algorithm can learn about any single input
record. The Gaussian mechanism is the canonical noise-addition primitive
used here.

\paragraph{DP-SGD.} Abadi et~al.~\cite{abadi2016dpsgd} adapted DP to
deep-learning training by clipping per-sample gradients before adding
calibrated Gaussian noise. The framework uses this exact construction at
the local-training step (\S\ref{sec:dpsgd-formal}). The two follow-ups
that drive accountant tightness are Mironov's Rényi differential
privacy~\cite{mironov2017rdp} (which gives the tightest known
$(\varepsilon, \delta)$ bounds for Gaussian DP-SGD) and the Privacy Random
Variable accountant~\cite{gopi2021prv} (slightly tighter still, supported
by Opacus as an alternative); the framework uses RDP by default.

\paragraph{DP-FedAvg.} McMahan et~al.~\cite{mcmahan2018dpfedavg} extended
DP-SGD to the federated setting with user-level privacy guarantees,
applying noise at the server-side aggregation step rather than the
client-side gradient step. Our construction places noise at the client
side (closer to standard DP-SGD), which gives \emph{record-level} DP
within each client; user-level DP would require either client-side
sampling at session boundaries or server-side noise. The trade-off is
discussed in Chapter~\ref{ch:disc}, \S7.4.

\paragraph{Opacus.} The Yousefpour et~al.\ Opacus
library~\cite{yousefpour2021opacus} is the framework's DP-SGD
implementation. Opacus contributes the per-sample-gradient hook
machinery, the RDP accountant, and the API for auto-calibrating $\sigma$
to a target $\varepsilon$. The framework's
\texttt{differential\_privacy.attach\_dp} is a thin wrapper around
\texttt{opacus.PrivacyEngine.make\_private\_with\_epsilon}.

\paragraph{Wei et~al.\ on DP-FL.} Wei et~al.~\cite{wei2020dpfl} provide
the most thorough empirical study of DP-FL utility-privacy tradeoffs to
date. They report that DP-FL on balanced benchmarks pays a small,
predictable utility cost; our work extends their analysis to imbalanced
healthcare data and finds a structurally different cost (class-prior
collapse, Chapter~\ref{ch:dp}).

\section{Secure aggregation}\label{sec:rw-secagg}

\paragraph{Bonawitz et~al.\ pairwise mask.} The secure-aggregation
protocol of Bonawitz et~al.~\cite{bonawitz2017secagg} is the framework's
direct ancestor. The protocol uses Diffie--Hellman key exchange to
derive pairwise seeds, expands the seeds via a PRG into mask vectors,
and uses Shamir secret-sharing of the seeds to tolerate up to $t$ client
dropouts per round. Each client masks its update; the masks cancel under
summation, so the server sees only the aggregate. Our prototype
implements the cancellation property (\S\ref{sec:secagg-formal}) but not
the dropout tolerance --- the simulation never drops a client.

\paragraph{Subsequent SecAgg variants.} A line of follow-up work
optimises SecAgg in different directions: SecAgg+~\cite{bell2020secaggplus}
reduces communication via a polylogarithmic communication graph;
LightSecAgg~\cite{so2022lightsecagg} eliminates the quadratic pairwise-mask
cost via single-mask broadcast and Lagrange-coded computing. These are
deployment-grade upgrades that the prototype does not need in order to
demonstrate the cancellation property.

\paragraph{TEE-based alternatives.} Trusted-execution-environment-based
schemes (Intel SGX, AMD SEV) provide aggregation confidentiality without
cryptographic masking, at the cost of a hardware trust assumption. We do
not use them because the framework is hardware-agnostic by design.

\section{Gradient compression}\label{sec:rw-compress}

\paragraph{Communication efficiency in FL.} Konečný
et~al.~\cite{konecny2016fedlearning} introduced the structured-update and
sketched-update families that compress FL traffic by two orders of
magnitude with bounded utility loss. Our top-$k$+int8 scheme is the
unstructured-sparsification + scalar-quantisation member of that family.

\paragraph{signSGD and 1-bit SGD.} Bernstein
et~al.~\cite{bernstein2018signsgd} push compression to its extreme:
transmit only the sign of each gradient component. signSGD trades
bandwidth for variance and has known convergence issues in non-IID FL,
which is why we use top-$k$ instead.

\paragraph{Top-$k$ sparsification.} Aji and
Heafield~\cite{aji2017topk} demonstrated that retaining only the
$\sim$1\% largest-magnitude gradient components per layer preserves
training quality on language modelling. We use $k = 0.1$ (10\%) which is
conservative for an imbalanced classification workload where rare-class
gradients are already small.

\paragraph{Compression with privacy.} Top-$k$ sparsification interacts
with DP-SGD in subtle ways: the noise that DP adds at the gradient step
becomes concentrated in a smaller index set when post-hoc sparsified,
amplifying its effect on the surviving entries. This is consistent with
our Finding 6.4 (\S\ref{sec:f6.3}): top-$k$ + int8 on top of DP-SGD pushes
the model back across the cliff on records 200--234, where DP-SGD alone
clears it.

\section{Privacy-preserving FL in healthcare}\label{sec:rw-health}

\paragraph{FL on ECG.} Kachuee et~al.~\cite{kachuee2018ecg} provide the
centralized baseline we reproduce; their reported test accuracy on
MIT-BIH is in the 0.97--0.99 range, consistent with the framework's
0.988 (\S\ref{sec:f4.1}). The published FL-on-ECG papers we surveyed
report similar headline accuracy under FedAvg or its variants but do
\emph{not} report majority-class baseline alongside top-1 accuracy ---
the editorial gap that motivates Finding 6.1 of this thesis.

\paragraph{FL on healthcare more broadly.} Rieke
et~al.~\cite{rieke2020future} survey the deployment landscape of
federated learning across medical imaging, EHR, and wearable health
data, and identify class imbalance as one of three open challenges (the
other two being client heterogeneity and protocol robustness). Our
class-prior collapse finding is a concrete instance of the imbalance
challenge they flag.

\paragraph{Regulatory framing.} HIPAA (Health Insurance Portability and
Accountability Act 1996, US Public Law 104-191) and GDPR (Regulation
(EU) 2016/679) are the two regulatory regimes that motivate the
thesis's data-minimisation framing. Federated learning is a natural
technical instantiation of the data-minimisation principle: training
data does not leave the device.

\section{Attacks: membership inference and gradient inversion}\label{sec:rw-attacks}

\paragraph{Membership inference.} Shokri
et~al.~\cite{shokri2017membership} demonstrated that an adversary with
black-box query access to a trained model can determine, with non-trivial
accuracy, whether a given record was in the training set. This is the
canonical attack against adversary A4 in Chapter~\ref{ch:threat}, and the
attack DP is built to defeat. The Yeom
et~al.~\cite{yeom2018membership} privacy-loss formalisation maps
membership-inference advantage to a function of $\varepsilon$, providing
a quantitative connection between the DP budget and attack success.

\paragraph{Gradient inversion.} Geiping
et~al.~\cite{geiping2020inverting} showed that an honest-but-curious
server can reconstruct individual training images with high fidelity from
the per-client gradient updates of a federated round, even under
multi-step FedAvg-style local training. The attack generalises beyond
images to other modalities and is the canonical attack against adversary
A1 in Chapter~\ref{ch:threat}. SecAgg is the canonical defence: under
SecAgg the server only ever sees the aggregate $\sum_i \Delta w_i$, and
gradient inversion against the aggregate is a strictly harder problem
than against a single client's update.

\paragraph{Property inference.} Beyond per-record attacks, property
inference~\cite{melis2019exploit} extracts aggregate statistics about the
training set (e.g.\ the gender ratio of contributing patients) from model
updates. The DP and SecAgg layers both partially mitigate property
inference; a complete defence would require formalising group-level DP,
which is future work.

\section{What is not covered}

The framework does not engage with three adjacent literatures that the
threat model deliberately excludes (Chapter~\ref{ch:threat}):
Byzantine-robust aggregation~\cite{blanchard2017krum, yin2018byzantine}
(median, trimmed mean, Krum), free-rider detection (clients that submit
zero or replay updates), and verifiable computation (cryptographic
proof that the server aggregated correctly). All three are listed in
\S\ref{sec:future} as natural extensions that compose cleanly on top of
the existing privacy stack.

\chapter{Threat model}\label{ch:threat}

The privacy stack of this framework is only as meaningful as the adversary
it is built against. This chapter enumerates the assets being protected, the
adversaries the framework defends against, the capabilities each is granted,
which component of the privacy stack mitigates each threat, and the trust
assumptions inherited from the deployment substrate. The threat model is the
contract that the design choices in Chapter~\ref{ch:method} are built to
honour, and the limitations called out in Chapter~\ref{ch:disc} are the
places where the contract is honoured only partially.

\section{Assets being protected}

Five categories of asset flow through the framework, each with its own
sensitivity profile (Table~\ref{tab:assets}). Raw waveforms and labels are
both protected health information (PHI) under HIPAA and personal data under
GDPR; per-client gradient updates leak training data via gradient-inversion
attacks even though they are nominally derived data; the aggregated global
model is the published artefact that membership-inference attacks operate on;
patient and device identifiers are quasi-identifiers that must not be
re-linkable to the local partitions.

\begin{table}[h]
  \centering
  \caption{Assets the framework protects and their sensitivity. The privacy
  stack is responsible for preventing leakage of every row except patient
  identity, which is handled at the deployment layer (no identifier ever
  leaves the device in our protocol).}
  \label{tab:assets}
  \begin{tabular}{ll}
    \toprule
    Asset & Sensitivity \\
    \midrule
    Raw ECG waveforms                       & PHI under HIPAA; highly sensitive \\
    Local labels (arrhythmia class)         & Diagnostic; highly sensitive \\
    Per-client model updates ($\Delta w_i$) & Leaks training data via gradient inversion \\
    Aggregated global model $w_t$           & Leaks via membership inference \\
    Patient / device identity               & Quasi-identifier; must not be re-linkable \\
    \bottomrule
  \end{tabular}
\end{table}

\section{Adversaries and capabilities}

Five adversary archetypes are considered. The first four are inside the
contract; the fifth is explicitly out of scope and motivates a future-work
item in Chapter~\ref{ch:disc}.

\subsection{A1 --- Honest-but-curious aggregation server}

\paragraph{Capability.} Sees every message that reaches the server: per-client
parameter updates, per-client metadata (number of training samples, training
loss, latency telemetry), broadcast parameters, and round-level scheduling
information.

\paragraph{Behaviour.} Follows the protocol faithfully. Does not deviate from
the aggregation rule, drop messages, or inject crafted messages.

\paragraph{Goal.} Infer training-data properties of any individual client
from its update --- e.g.\ reconstruct individual training records from
$\Delta w_i$ via gradient-inversion attacks, or learn whether a particular
patient's beats appear in client $i$'s partition.

\paragraph{Mitigation.} Pairwise-mask secure aggregation
(Chapter~\ref{ch:method}, \S\ref{sec:method-strategy}). Under SecAgg, the
server's view of every per-round message is a uniformly random vector;
only the sum across all clients de-randomises into the true aggregate of
deltas. The server therefore has access to $\sum_i \Delta w_i$ but not to
any individual $\Delta w_i$.

\subsection{A2 --- Curious-client coalition (size $t < N$)}

\paragraph{Capability.} A subset of clients colludes among themselves. They
pool their private masks (or shared seeds in the prototype), their broadcast
parameters, and the metadata they each receive from the server.

\paragraph{Behaviour.} Trains honestly so as not to be detected by simple
correctness checks (the framework's SecAgg is not currently
correctness-verifiable from the server's side; that is future work).

\paragraph{Goal.} Unmask a non-coalition client's update.

\paragraph{Mitigation.} Pairwise masks between every unordered client pair
$\{i,j\}$ mean that a coalition of size $t$ learns nothing about non-coalition
updates as long as at least one mask between an honest pair survives the
collusion. With $N = 4$ and pairwise masks the basic scheme withstands a
coalition of up to $t = N - 2 = 2$ without compromise; with $t = N - 1 = 3$
the scheme leaks the remaining update because the only mask the honest
client uses is one whose seed is held by a coalition member.

\subsection{A3 --- Network eavesdropper}

\paragraph{Capability.} Passive observation of the wire between client and
server (e.g.\ an on-path attacker on Wi-Fi, an ISP, or a nation-state
adversary at the network layer).

\paragraph{Goal.} Same as A1, but with the additional restriction that the
adversary does not control the server's storage or compute.

\paragraph{Mitigation.} Transport-layer encryption (TLS, assumed; out of
scope for the prototype but standard in any production deployment). The
masked SecAgg payload is also a uniformly random vector to a network
observer, providing defence-in-depth: even if TLS is subsequently broken,
A3's view is no better than A1's, which SecAgg already addresses.

\subsection{A4 --- External adversary holding the released global model}

\paragraph{Capability.} Inspects, queries, or fine-tunes the released global
model weights $w_T$. Does not see per-round messages. Has access to
auxiliary information about all but one of the training records (the
worst-case threat model for differential privacy).

\paragraph{Goal.} Membership inference: ``was patient $P$'s data in the
training set?'' or its quantitative variant, ``what fraction of patients in
this hospital appear in the training set?''. Reconstruction: recover
features of the training distribution beyond what is published.

\paragraph{Mitigation.} \textbf{Differential privacy} with budget
$(\varepsilon, \delta)$ at the per-record level. With $\varepsilon = 5$ and
$\delta = 10^{-5}$ (the framework's default targets) the worst-case posterior
shift on any single-record membership question is at most $e^{\varepsilon}
\approx 148$, and empirical membership-inference accuracy is typically much
closer to the $50\%$ random-guess baseline. The cost of this guarantee on
imbalanced healthcare data is the central finding of Chapter~\ref{ch:dp}.

\subsection{A5 --- Malicious (Byzantine) client}

\paragraph{Capability.} Deviates arbitrarily from the protocol: submits
crafted updates designed to poison the global model, install a backdoor,
or break SecAgg by sending mismatched mask contributions.

\paragraph{Goal.} Either degrade the global model's utility, install a
trigger pattern that produces attacker-chosen outputs at inference time,
or extract information from honest clients by exploiting protocol
deviations.

\paragraph{Out of scope.} The framework's threat model does not include A5.
Robust aggregation rules (Krum, coordinate-wise median, trimmed mean,
multi-Krum) and protocol-level integrity checks would be the natural
defences; both are listed as future work in Chapter~\ref{ch:disc},
\S\ref{sec:future}. The simulation never produces an A5: the four clients
are honest by construction.

\section{Attack-surface coverage matrix}

Table~\ref{tab:attack-surface} maps the five canonical attacks against
federated learning to which adversary mounts them and which layer of the
privacy stack mitigates each one. The matrix is the most compact statement
of the contract: every cell either references a mitigation or is explicitly
out of scope. There are no implicit guarantees.

\begin{table}[h]
  \centering
  \caption{Attack-surface coverage. Cells marked \emph{SecAgg} are mitigated
  by the pairwise-mask layer, cells marked \emph{DP} by differential privacy,
  cells marked \emph{TLS+SecAgg} by both, and cells marked \emph{OOS} are
  out of scope (Byzantine adversary A5).}
  \label{tab:attack-surface}
  \begin{tabular}{lccccc}
    \toprule
    Attack & A1 & A2 & A3 & A4 & A5 \\
    \midrule
    Gradient inversion                    & SecAgg     & SecAgg     & TLS+SecAgg & ---  & OOS \\
    Per-update membership inference       & SecAgg     & SecAgg     & TLS+SecAgg & ---  & OOS \\
    Global-model membership inference     & DP         & DP         & ---        & DP   & OOS \\
    Reconstruction from compressed update & SecAgg     & SecAgg     & TLS+SecAgg & ---  & OOS \\
    Model poisoning / backdoor            & ---        & ---        & ---        & ---  & OOS \\
    \bottomrule
  \end{tabular}
\end{table}

\paragraph{Why compression does not weaken the matrix.}
Top-$k$ + int8 compression sits between local DP-SGD and SecAgg masking on
the client (Chapter~\ref{ch:method}, \S\ref{sec:method-strategy}). Because
SecAgg is applied \emph{after} compression, the wire payload is still a
uniformly random vector to A1/A2/A3, and the compressed update is never
visible in the clear. DP's post-processing immunity ensures that compression
does not weaken A4's bound on the released global model. Compression is a
bandwidth optimisation, not a privacy defence; it is included in the matrix
only to make explicit that it does not subtract from the matrix either.

\section{Trust assumptions}\label{sec:trust}

The framework's guarantees are conditional on four trust assumptions
inherited from the deployment substrate:

\begin{enumerate}[leftmargin=*]
  \item \textbf{TLS is in place between every client and the server} ---
  standard mTLS in any production deployment. The simulation does not
  exercise this; in a real IoT deployment it is non-negotiable.
  \item \textbf{The pseudo-random generator used to derive masks is
  cryptographically secure.} The prototype uses Python's \texttt{secrets}
  module for seed generation and a deterministic NumPy generator from that
  seed for mask material; the latter is sufficient for simulation but a
  production deployment would seed an authenticated stream cipher (e.g.\
  AES-CTR) instead.
  \item \textbf{No more than $N - 2$ clients collude.} The basic
  pairwise-mask scheme has no Shamir secret-sharing fallback (a deliberate
  prototype simplification, called out in
  \texttt{src/privacy/secure\_agg.py}), so a coalition of $N - 1$ clients
  reconstructs the remaining client's update.
  \item \textbf{Hardware-level side channels (timing, power, EM) on the IoT
  device are out of scope.} A real wearable deployment would also need
  hardware countermeasures; without them, an attacker with physical access
  to a device can extract the local model or training data regardless of
  what the FL protocol does.
\end{enumerate}

\section{What is explicitly out of scope}

Beyond A5, four classes of guarantee are explicitly not provided:

\begin{itemize}[leftmargin=*]
  \item \textbf{Byzantine robustness / poisoning defences.} Future work
  (\S\ref{sec:future}); the natural extension is robust aggregation
  (Krum, median, trimmed mean) composed with the existing privacy stack.
  \item \textbf{Hardware-level side channels.} Out of scope at the protocol
  layer; addressed at the deployment layer with hardware countermeasures.
  \item \textbf{Verifiable computation} (cryptographic proof that the
  server aggregated correctly). The framework trusts that the server runs
  the strategy code it claims to run; deployments in adversarial settings
  would add zk-SNARK or trusted-execution-environment-based attestation.
  \item \textbf{Free-rider detection.} A client that submits zero-update
  vectors round after round contributes nothing to the global model but
  receives the benefit of others' updates; the framework does not detect
  or penalise this. Reputation and contribution-tracking systems are the
  standard defence; we do not implement either.
\end{itemize}

These are flagged again in the conclusions chapter as natural extensions
that fit cleanly on top of the existing privacy stack.

\chapter{Method}\label{ch:method}

This chapter describes the implementation that produces every number in the
thesis. All code lives under \texttt{src/} and \texttt{experiments/}; YAML
configs in \texttt{configs/} are the canonical record of every run. The
repository is at
\url{https://github.com/mapeke/privacy-preserving-fl-healthcare-iot}.

\section{Unified client interface (\texttt{src/client/flower\_client.py})}

\texttt{FlowerHealthcareClient} subclasses \texttt{flwr.client.NumPyClient}
with three methods. \texttt{fit} performs local training (optionally
DP-SGD), computes the parameter delta against the broadcast parameters,
optionally compresses it, optionally adds the SecAgg pairwise masks, and
returns reconstructed parameters (broadcast~+~processed~delta) so vanilla
Flower strategies remain compatible. \texttt{evaluate} reports per-client
test loss and accuracy. \texttt{get\_parameters} serialises the model state
into the wire format. The class takes a single \texttt{ClientConfig} dataclass
with mutually-orthogonal fields \texttt{dp\_config}, \texttt{secagg\_session},
\texttt{compression\_method}, each of which is \texttt{None} when the
corresponding defense is off.

\section{Data layer (\texttt{src/data/})}\label{sec:method-data}

\texttt{mitbih\_loader.py} downloads MIT-BIH records via
\texttt{wfdb.dl\_database} into \texttt{data/cache/}, segments each record
into 360-sample beat windows centered on R-peak annotations, applies
per-record z-score normalisation, and maps annotation symbols to AAMI labels.
A synthetic-ECG fallback exists for offline development; both real-data and
synthetic-data runs go through the same downstream pipeline.

\texttt{partitioner.py} distributes beats across $N$ clients via Dirichlet
sampling per class with concentration $\alpha$ (default 0.5; lower
$\Rightarrow$ more skewed). For our 4-client setup the realised partition
sizes are $(1907, 8621, 6947, 2840)$ --- non-IID by construction.

\texttt{preprocessing.py} provides a Butterworth bandpass filter (0.5--40
Hz) for baseline-wander and powerline rejection, a stratified train/test
split, and the \texttt{TorchEcgDataset} adapter.

\section{Models (\texttt{src/models/})}

\texttt{EcgCNN} is a 3-block 1-D CNN with GroupNorm
(used in place of BatchNorm for Opacus compatibility), 64 final channels,
adaptive average pooling, and a 32-unit MLP head; $\sim$50k parameters.
\texttt{LightweightEcgCNN} is a depthwise-separable variant at $\sim$5k
parameters for resource-constrained edge inference. All experiments in this
thesis use \texttt{EcgCNN}.

\paragraph{ReLU non-inplace.}
\texttt{nn.ReLU(inplace=True)} is incompatible with Opacus's per-sample
gradient hooks (the in-place op writes through a backward-hook view).
The exception is swallowed by the Ray worker process in Flower's
simulation backend, producing silently-flat training. Switching to
non-inplace ReLU fixed this; see commit \texttt{40dd3f4}.

\section{Differential privacy (\texttt{src/privacy/differential\_privacy.py})}

\texttt{DPConfig} is the user-facing dataclass with
\texttt{target\_epsilon}, \texttt{target\_delta}, \texttt{max\_grad\_norm}
(the per-sample clip threshold $C$), optional \texttt{noise\_multiplier}, and
\texttt{accountant} ($\in\{$\texttt{rdp}, \texttt{gdp}, \texttt{prv}$\}$).
\texttt{attach\_dp(model, optimizer, dataloader, config, epochs)} wraps the
optimizer and dataloader via \texttt{opacus.PrivacyEngine}. When
\texttt{noise\_multiplier} is left unset Opacus auto-calibrates $\sigma$ to
hit \texttt{target\_epsilon} over the requested epochs; otherwise the
user-supplied $\sigma$ is used. The returned handle exposes
\texttt{get\_epsilon(steps)} for per-round budget reporting.

\section{Secure aggregation (\texttt{src/privacy/secure\_agg.py})}

\texttt{setup\_session(num\_clients, seed\_size\_bytes=32, rng\_seed=...)}
generates one fresh secret seed per unordered client pair, stored in a
\texttt{SecAggSession} object distributed by the experiment orchestrator.
\texttt{derive\_mask(seed, shape)} expands a pair seed into a Gaussian
mask vector via NumPy's PCG64. \texttt{mask\_update(update, client\_id,
session)} adds $+m_{ij}$ for $j > i$ and $-m_{ij}$ for $j < i$ to the
client's update, so that summing all clients' masked updates cancels every
mask exactly. \texttt{aggregate\_masked\_updates} performs the server-side
sum on the strategy side
(\S\ref{sec:method-strategy}).

\paragraph{Limitations explicitly accepted.}
Pre-shared seeds rather than Diffie--Hellman; no Shamir secret-sharing
(no dropout tolerance); float32 arithmetic rather than modular integer over
a finite field. These are appropriate for a diploma simulation that wants
to demonstrate the principle and measure end-to-end overhead, and they are
flagged in the module docstring.

\section{Compression (\texttt{src/privacy/compression.py})}

\texttt{compress(update, method, top\_k\_ratio, quant\_bits)} dispatches over
four modes: \texttt{none}, \texttt{top\_k\_only}, \texttt{quantize\_only},
\texttt{top\_k\_quantize}. \texttt{top\_k\_sparsify} keeps the largest-magnitude
$\lceil k|u|\rceil$ entries via $O(n)$ argpartition.
\texttt{quantize\_int8} performs min--max quantisation with a documented
\texttt{(scale, zero\_point)} pair; \texttt{decompress} reconstructs a dense
update array (zeros where sparsified). The wire payload is a
\texttt{CompressedUpdate} dataclass with \texttt{indices}, \texttt{values},
\texttt{original\_shape}, and a \texttt{num\_bytes()} method for bandwidth
accounting.

\section{Server strategy (\texttt{src/server/strategies.py})}\label{sec:method-strategy}

\texttt{PrivacyAwareFedAvg} subclasses Flower's \texttt{FedAvg}. With SecAgg
disabled it falls through to standard weighted averaging. With SecAgg
enabled, \texttt{configure\_fit} caches the broadcast parameters;
\texttt{aggregate\_fit} recovers per-client deltas by subtracting them from
each client's returned parameters, sums the deltas (where masks cancel by
construction), divides by the number of participating clients, and applies
the result on top of the cached broadcast. Robust scalar-metric aggregation
(sample-weighted mean for numerics, first-occurrence for non-numerics)
covers the case where Flower's default aggregator chokes on
heterogeneous metric dictionaries.

\section{Experiment infrastructure}

Five entry points, each a thin CLI wrapper around
\texttt{experiments/\_common.py::build\_experiment} +
\texttt{run\_federated\_experiment}:

\begin{itemize}[leftmargin=*]
  \item \texttt{run\_centralized.py} --- pooled-data baseline.
  \item \texttt{run\_federated.py} --- vanilla FedAvg.
  \item \texttt{run\_federated\_dp.py} --- FedAvg + DP-SGD.
  \item \texttt{run\_federated\_secagg.py} --- FedAvg + SecAgg.
  \item \texttt{run\_full\_pipeline.py} --- DP + SecAgg + compression.
\end{itemize}

Each writes \texttt{results/<name>/\{run.log, resolved\_config.yaml,
history.json\}} so every reported number can be traced back to its config.

\section{Tests}\label{sec:method-tests}

42 pytest unit tests across seven modules; the most consequential is
\texttt{test\_secagg.py::test\_masks\_cancel\_under\_aggregation}, which
constructs four random updates, masks them, sums the masked vectors,
and asserts that the result equals the unmasked sum to within
$10^{-4}$ in $\ell_2$. This is contribution C2: the cryptographic
correctness argument as a code-level invariant.

The test suite passes in 4.0s on the developer machine. \texttt{Opacus} is
imported lazily inside \texttt{attach\_dp} so the tests run without it
installed; the test that asserts that lazy-import property is itself an
invariant.

\chapter{Results I --- Centralized and federated baselines}\label{ch:detection}

Setup: 12 MIT-BIH records (Table~\ref{tab:class-counts}), 4 clients,
Dirichlet $\alpha = 0.5$ partitions of sizes $(1907, 8621, 6947, 2840)$,
80/20 stratified train/test split.
Configs: \texttt{configs/default.yaml}.
Hardware: single NVIDIA RTX (mobile) GPU, CUDA 12.1, PyTorch 2.4.1.
Seed 42.

\section{Detection grid}

Table~\ref{tab:detection-grid} reports the headline accuracy for every
configuration we run in this thesis. The remainder of this chapter establishes
the top two rows; Chapters~\ref{ch:secagg} and \ref{ch:dp} establish the rest.

\begin{table}[h]
  \centering
  \caption{Final test accuracy across all five experimental configurations.
  Federation is free; SecAgg is free; DP collapses to the majority-class
  prevalence; compression on top of DP does not move the needle further.
  Numbers are single-seed at seed 42 to match the thesis policy of
  ``every number tied to one config row.''}
  \label{tab:detection-grid}
  \begin{tabular}{lrrr}
    \toprule
    Configuration & Test acc & Macro-F1 & Final $\varepsilon$ \\
    \midrule
    Centralized (5 epochs)              & \textbf{0.9878} & 0.36 & --- \\
    FedAvg, 4 clients, 10 rounds        & \textbf{0.9880} & ---  & --- \\
    FedAvg + SecAgg, 10 rounds          & 0.9860          & ---  & --- \\
    FedAvg + DP, 20 rounds, target $\varepsilon = 5$
                                        & 0.9711          & ---  & 4.9955 \\
    Full pipeline (DP+SecAgg+Comp), 30 rounds, target $\varepsilon = 8$
                                        & 0.9711          & ---  & 7.9964 \\
    \midrule
    Class-\emph{N} prevalence (= predict-all-N)
                                        & 0.9716          & 0.20 & --- \\
    \bottomrule
  \end{tabular}
\end{table}

\section{Centralized baseline}

5 epochs of plain SGD on the pooled training set (20{,}315 beats),
$\eta = 0.01$, momentum 0.9, weight decay $10^{-4}$, batch size 32.
Final test accuracy 0.9878, macro-F1 0.3627. Per-class F1:
$\{N{=}0.99,\ V{=}0.82,\ S{=}0.00,\ F{=}0.00,\ Q{=}0.00\}$. The model
learns the two prevalent classes (N, V) and ignores the others. Macro-F1 is
$\sim$0.36 because three of the five classes contribute zero, a structural
property of the workload.

\section{Vanilla FedAvg}\label{sec:f4.1}

10 rounds, 4 clients per round, 2 local epochs, batch 32, $\eta = 0.01$.
Round-by-round test accuracy: $\{0.9711, 0.9711, 0.9758, 0.9840, 0.9840,
0.9724, 0.9831, 0.9864, 0.9862, 0.9880\}$.

\paragraph{Finding 4.1 --- Federation is free on this workload.}
Final FedAvg accuracy 0.9880 vs.\ centralized 0.9878.
The federated decomposition introduces zero accuracy cost up to seed
variance. No special tuning, no proximal regularizer~\cite{li2020fedprox}, no
server-momentum --- plain FedAvg suffices when client data is similar enough
in distribution and the model is not over-parameterised.

This is the empirical baseline against which Chapters~\ref{ch:secagg} and
\ref{ch:dp} measure the cost of the privacy stack. Any utility loss observed
with SecAgg or DP enabled is attributable to the privacy mechanism, not to
the federation itself.

\section{Reproducibility check vs.\ published FL ECG results}

Single-seed run; we did not run multi-seed for the headline figure due to
GPU-time budget (each FedAvg run is $\sim$3 minutes on our hardware, but
DP and full-pipeline runs are 30 minutes each; multi-seed across all
five configurations is a future-work item, \S\ref{sec:future}). Published FL
ECG papers report top-1 accuracy in the 0.97--0.99 range on similar MIT-BIH
subsets~\cite{kachuee2018ecg}; we are within that range. The thesis's
contribution does not depend on matching any published absolute --- it
depends on the relative shape across the privacy stack rows.

\chapter{Results II --- Secure aggregation}\label{ch:secagg}

Configs: \texttt{configs/experiment\_secagg.yaml}. SecAgg enabled,
DP off, compression off. 10 rounds. Otherwise identical to
Chapter~\ref{ch:detection}.

\section{Mask cancellation as a code-level invariant}\label{sec:f5.1}

\paragraph{Finding 5.1 --- Cryptographic correctness is mechanically
verifiable.} The pairwise-mask correctness argument
\cite{bonawitz2017secagg} is reduced to one assertion:
\begin{quote}\itshape
  for $N=4$ random updates of shape $(10,)$,
  $\sum_{i=1}^{N}\mathrm{mask\_update}(u_i, i, \mathrm{session}) =
   \sum_{i=1}^{N} u_i$ to within $10^{-4}$ $\ell_2$ error.
\end{quote}
This lives in
\texttt{tests/test\_secagg.py::test\_masks\_cancel\_under\_aggregation} and
runs in CI alongside the rest of the pytest suite. The same test verifies
the dual property: each individual masked update is \emph{not} close to the
unmasked update (the masks are non-trivial; the cancellation only happens
at the sum level).

The contribution here is not the test itself --- it is one screenful of
NumPy --- but the discipline of writing the protocol's correctness argument
as a checked invariant rather than as prose. A future change that breaks the
protocol (e.g., an off-by-one in the $i < j$ vs.\ $i > j$ sign condition)
fails the test loudly rather than silently producing biased aggregates.

\section{End-to-end SecAgg run}\label{sec:f5.2}

Round-by-round test accuracy: $\{0.9711, 0.9711, 0.9711, 0.9776, 0.9837,
0.9840, 0.9844, 0.9848, 0.9846, 0.9860\}$. Final 0.9860.

\paragraph{Finding 5.2 --- SecAgg overhead is below run-to-run noise.}
Vanilla FedAvg final 0.9880 vs.\ FedAvg + SecAgg final 0.9860 ---
a 0.20\,pp gap.
Per-class F1, train accuracy, and convergence trajectory all overlap
with vanilla FedAvg's within a few hundredths.
The $0.20\,\mathrm{pp}$ gap is consistent with float32 accumulation noise
in the masking; it would shrink under integer modular arithmetic
(future work, \S\ref{sec:future}).

The combination of \S\ref{sec:f5.1} and this finding is the case for
SecAgg as a default in any FL deployment: the cryptographic
correctness is a one-line CI check; the utility cost is below the noise
floor; and the threat model that SecAgg addresses (honest-but-curious
server with read access to the inbox) is realistic for any cloud-hosted
aggregation server.

\section{What SecAgg does not protect against}

SecAgg blinds the server to per-client updates but does not blind it
to the \emph{aggregate}. An adversary holding the released global
model $w_T$ can still mount membership-inference attacks against
training data; that threat is what differential privacy
(Chapter~\ref{ch:dp}) addresses. The two layers are complementary,
and the thesis argues for running both, modulo the utility cost
which Chapter~\ref{ch:dp} now measures.

\chapter{Results III --- Differential privacy and the full pipeline}\label{ch:dp}

Configs: \texttt{configs/experiment\_dp.yaml} and
\texttt{configs/experiment\_full.yaml}. DP enabled (Opacus PrivacyEngine,
RDP accountant, \texttt{max\_grad\_norm = 1.0}). Otherwise identical to
Chapter~\ref{ch:detection}.

\section{DP-SGD run}

20 rounds, 1 local epoch per round, batch size 64, $\eta = 0.005$,
target $\varepsilon = 5.0$, target $\delta = 10^{-5}$.
Opacus auto-calibrates $\sigma$ to hit
$\varepsilon = 4.9955$ (well within the target).

Round-by-round test accuracy:
$\{0.9711, 0.9711, 0.9711, 0.9711, 0.9711, \ldots, 0.9711, 0.9711, 0.9711\}$.
The trajectory is exactly flat to four decimal places across all
20 rounds.

Train metrics tell the inner story:
train loss falls 0.5370 $\to$ 0.2233; train accuracy rises 0.9147 $\to$
0.9712. The model is doing \emph{something} during training; that
something just happens to be ``learn to predict N and only N.''

\section{Class-prior collapse}\label{sec:f6.1}

\paragraph{Finding 6.1 --- DP-SGD collapses to majority-class prediction.}
Class \emph{N} prevalence in the test split: $4{,}934 / 5{,}078 = 0.9716$.
Final DP-SGD test accuracy: $0.9711$. The two numbers agree to four
decimals. The model is, with statistical certainty, the majority-class
predictor.

This is not a learning-rate or noise-multiplier failure --- the train
accuracy converges to the same 0.9712, and the spent $\varepsilon$ is
within the target. It is a workload-mismatch failure: DP-SGD's per-step
Gaussian noise drowns out the gradient signal of the rare classes, which
contribute $\sim$2.9\,\% of the data. Under sufficient noise, the
empirical-risk-minimising policy on the training set is exactly:
predict the majority. Opacus is doing what it advertises; the workload
asks more of it than the budget will pay for.

\paragraph{Why this matters editorially.}
A casual reader of Table~\ref{tab:detection-grid} sees ``0.9711''
beside ``DP $\varepsilon = 5$'' and concludes that DP costs the
framework about 1.7\,pp of accuracy --- a number that would be
publishable under the headline-accuracy convention dominant in the
FL-privacy literature. The honest reading is that DP costs the model's
ability to detect any rare class at all on this data. The 0.9711 number
hides the failure precisely because the metric is overall accuracy on a
heavily imbalanced label space.

The framework therefore reports macro-F1 alongside accuracy in
\texttt{src/utils/metrics.py::compute\_classification\_report}; the
centralized run's macro-F1 of 0.36 (Chapter~\ref{ch:detection}) drops to
$\le 0.20$ on the DP run, which would have made the collapse visible
if multi-seed F1 had been logged across rounds. We did not log it for
the federated runs; we will (\S\ref{sec:future}).

\section{The full pipeline run}

30 rounds, 1 local epoch, batch 64, $\eta = 0.005$, target
$\varepsilon = 8.0$, target $\delta = 10^{-5}$, SecAgg enabled,
top-$k$ + int8 compression with $k = 0.1$.
Opacus auto-calibrates $\sigma$ for the looser $\varepsilon$, hits
$\varepsilon = 7.9964$.

Round-by-round test accuracy: 0.9711 across all 30 rounds.

\paragraph{Finding 6.2 --- Compression on top of DP does not degrade further.}
DP-only at $\varepsilon = 5$ converges to 0.9711. Full pipeline (DP +
SecAgg + compression) at $\varepsilon = 8$ converges to 0.9711.
The model is already at the floor that the workload's class prior
imposes; there is nothing left to degrade by adding more layers.
Loosening the privacy budget from $\varepsilon = 5$ to $\varepsilon = 8$
does not move the model off the floor either, which would otherwise be
the natural follow-up question.

\paragraph{Compositional implication.}
For an operator constrained to deploy DP at $\varepsilon \le 8$ on a
class-imbalanced workload, adding SecAgg and compression is essentially
free: they do not make the already-collapsed model any worse. This
inverts the usual ``every layer costs something'' framing for the
specific case where the bottom layer has already maxed out the
utility cost.

\section{Negotiating the cliff}\label{sec:negotiating}

The thesis does not propose a fix; the experimental contribution is to
make the cliff visible. Three directions a follow-up paper would take:

\begin{enumerate}[leftmargin=*]
  \item \textbf{Class-weighted DP-SGD.} Apply per-class loss weights so
  that the gradient signal from the rare classes is amplified before
  clipping. Risks: changes the privacy accounting because the
  $\ell_2$ sensitivity of the loss is no longer 1.
  \item \textbf{Per-class privacy budgets.} Decouple $\varepsilon$ for
  majority and minority classes, e.g.\ $\varepsilon_N = 1$,
  $\varepsilon_{S, V, F, Q} = 10$. Mathematically clean but practically
  awkward to deploy --- the privacy guarantee is then conditional on the
  target class, which a deployment operator may find harder to reason
  about than a single $\varepsilon$.
  \item \textbf{Different MIT-BIH records.} Records 200--234 contain
  more arrhythmias and reduce the imbalance from 97.1\,\% N to
  82.1\,\% N. We re-ran all five configurations on this slice
  (\S\ref{sec:f6.3} below) without retuning, both to test whether the
  cliff is dataset-conditional and to give the rare-class classifier a
  fairer chance. The deployment workload (a wearable monitoring an
  asymptomatic patient) is closer to the 100--115 imbalance than to
  200--234, so the primary slice remains the more representative one.
\end{enumerate}

\section{Follow-up: records 200--234}\label{sec:f6.3}

To test whether the DP class-prior collapse (\S\ref{sec:f6.1}) is a
property of DP-SGD itself or of the workload's extreme imbalance, we
re-ran all five experimental configurations on records 200--234 ---
a 25-record slice with 61{,}818 beats and substantially more
arrhythmia variety. No hyperparameters were retuned; configs are
\texttt{configs/followup\_records200/\{default, experiment\_dp,
experiment\_secagg, experiment\_full\}.yaml}.
Table~\ref{tab:class-balance-200} contrasts the two slices' class
distributions; Table~\ref{tab:followup-grid} reports final accuracies
side by side with the original Table~\ref{tab:detection-grid}.

\begin{table}[h]
  \centering
  \caption{AAMI class distribution on records 100--115 (the primary
  slice, 25{,}393 beats) and records 200--234 (the follow-up,
  61{,}818 beats). The follow-up reduces the majority-class share
  from 97.1\,\% to 82.1\,\% and gives all four rare classes between
  1\,\% and 10\,\% of the data each.}
  \label{tab:class-balance-200}
  \begin{tabular}{lrr}
    \toprule
    Class & 100--115 share & 200--234 share \\
    \midrule
    \emph{N} (normal)              & 97.10\,\% & 82.08\,\% \\
    \emph{S} (supraventricular)    &  0.24\,\% &  4.19\,\% \\
    \emph{V} (ventricular ectopic) &  2.60\,\% &  9.53\,\% \\
    \emph{F} (fusion)              &  0.03\,\% &  1.28\,\% \\
    \emph{Q} (unknown / paced)     &  0.03\,\% &  2.93\,\% \\
    \bottomrule
  \end{tabular}
\end{table}

\begin{table}[h]
  \centering
  \caption{Same five experiments, records 100--115 vs.\ records
  200--234, same hyperparameters in both columns. The majority-class
  baseline (predict-all-N) is 97.16\,\% on 100--115 and 82.08\,\% on
  200--234; the right-hand column is the gap from that baseline on
  the follow-up slice.}
  \label{tab:followup-grid}
  \begin{tabular}{lrrr}
    \toprule
    Configuration & 100--115 acc & 200--234 acc & vs.\ majority \\
    \midrule
    Centralized (5 epochs)              & 0.9878 & \textbf{0.9389} & $+11.81$\,pp \\
    FedAvg (10 rounds)                  & 0.9880 & \textbf{0.9636} & $+14.28$\,pp \\
    FedAvg + SecAgg (10 rounds)         & 0.9860 & \textbf{0.9558} & $+13.50$\,pp \\
    FedAvg + DP, $\varepsilon \approx 5$ (20 rounds)
                                        & 0.9711 & \textbf{0.8845} & $+6.37$\,pp \\
    Full pipeline, $\varepsilon \approx 8$ (30 rounds)
                                        & 0.9711 & \textbf{0.8176} & $-0.32$\,pp \\
    \bottomrule
  \end{tabular}
\end{table}

\paragraph{Finding 6.3 --- The DP class-prior collapse is
dataset-conditional.}
On records 200--234, DP-SGD at $\varepsilon \approx 5$ converges to
$0.8845$ test accuracy, $+6.37$\,pp above the $0.8208$ majority-class
baseline. The model has \emph{not} collapsed: it is distinguishing
rare-class beats. Compare to the original 100--115 result, where
DP-SGD's $0.9711$ matched class-\emph{N} prevalence to four decimals.
Finding 6.1 (\S\ref{sec:f6.1}) is therefore not a property of DP-SGD
but a property of \emph{extreme} class imbalance composed with DP's
per-step noise. The cliff exists; its location is workload-dependent.

\paragraph{Finding 6.4 --- The full pipeline collapses on moderate
imbalance too.}
The full-pipeline run (DP + SecAgg + top-$k$ + int8 at $\varepsilon
\approx 8$) on records 200--234 converges to $0.8176$, $0.32$\,pp
\emph{below} the $0.8208$ majority-class baseline. Round-by-round
accuracy oscillates between $\sim$65\,\% and $\sim$86\,\% over the
30 rounds; the run does not stabilise. Top-$k$ + int8 compression on
top of DP-SGD's already-noisy gradients appears to push the regime
back across the cliff that DP-SGD alone clears on this dataset.
Loosening the budget from $\varepsilon = 5$ (DP-only) to
$\varepsilon = 8$ (full pipeline) is not enough to compensate ---
the additional noise floor is the lossy compression, not the
DP budget.

\paragraph{What the follow-up changes editorially.}
The Chapter~\ref{ch:dp} framing ``DP collapses on imbalanced
healthcare data'' is true on records 100--115 but reads as a
generalisation. The honest generalisation is: \emph{DP-SGD's
signal-to-noise ratio degrades sharply when the rare-class share
approaches the per-step noise floor; the cliff is observable but its
location depends on both the budget and the workload imbalance.} The
full pipeline crosses the cliff on both slices, which is the more
generalisable finding the framework supports.

\paragraph{What the follow-up does not change.}
We do not retune hyperparameters per dataset --- the contribution is
the relative shape across the privacy stack, not absolute accuracy on
either slice. We do not claim records 200--234 is the ``right'' slice;
the deployment workload (a wearable monitoring an asymptomatic
patient) is closer to the 100--115 imbalance. Both slices are useful:
100--115 surfaces the cliff at the dataset extreme, 200--234 shows
where on the class-balance axis it sits.

\chapter{Discussion}\label{ch:disc}

\section{What the three findings add up to}

\begin{itemize}[leftmargin=*]
  \item \textbf{Federation (\S\ref{sec:f4.1}).} Free on this workload.
  \item \textbf{Secure aggregation (\S\ref{sec:f5.1}, \S\ref{sec:f5.2}).}
  Free up to fp32 mask-cancellation noise; cryptographic correctness is
  a one-line CI invariant.
  \item \textbf{Differential privacy (\S\ref{sec:f6.1}).} The
  privacy-utility cliff is real and is the central pitfall of DP-SGD on
  imbalanced healthcare data. The DP run's headline accuracy looks
  reasonable (0.9711) precisely because the metric matches the cliff:
  predicting the majority class scores 0.9716.
  \item \textbf{Dataset dependence (\S\ref{sec:f6.3}).} Finding 3 is
  conditional on \emph{extreme} class imbalance: on records 200--234
  (82\,\% N rather than 97\,\% N) DP-SGD alone retains a $+6.37$\,pp
  rare-class signal, while the full pipeline still collapses. The
  cliff is a property of the workload composed with the privacy stack,
  not of DP-SGD per se; the framework is what makes that property
  visible.
\end{itemize}

The thesis-level recommendation: any privacy-preserving FL paper for
imbalanced healthcare data should report majority-class baseline alongside
overall accuracy and macro-F1 alongside top-1; the privacy/utility tradeoff
graph is misleading without them. Our framework contributes both the
mechanism that produces the data and the protocol that surfaces the cliff.

\section{Likely reviewer objections}

\paragraph{``You only ran each configuration at one seed.''}
True. Vanilla FedAvg, SecAgg, and DP are seed-stable on this dataset to
within roughly 0.5\,pp accuracy according to spot-checks during
development. The DP collapse is structural rather than seed-dependent
--- it is the model converging to the majority-class predictor, which
the same training process produces from any random initialisation. Multi-seed
across all configurations is the first item of \S\ref{sec:future}.

\paragraph{``You only used 12 records.''}
The primary experiments use records 100--115 (the
\texttt{wfdb.dl\_database} default). \S\ref{sec:f6.3} re-runs all
five experiments on records 200--234 (25 records, 61{,}818 beats,
82.1\,\% N) without retuning. The qualitative finding holds with one
refinement: DP-SGD alone clears the cliff on records 200--234 but the
full pipeline still collapses. The 12-record default is enough to
\emph{surface} the cliff; the 25-record follow-up is what places it
on the class-balance axis.

\paragraph{``You did not compare to FedProx or SCAFFOLD.''}
The contribution is layering existing privacy primitives on top of
FedAvg, not optimising the FL aggregation itself. FedAvg achieves
0.988 on this workload --- approximately the centralized baseline.
There is no headroom for a stronger FL aggregator to demonstrate on
this data. A different workload (more clients, more heterogeneity)
would justify the comparison.

\paragraph{``Your SecAgg is not Bonawitz et~al.''}
Acknowledged. The prototype uses pre-shared seeds rather than
Diffie--Hellman, no Shamir secret-sharing for dropout tolerance,
and float32 arithmetic. These limitations are listed in the
module docstring, in
\texttt{docs/privacy-analysis.md\,\S2}, and at the start of
\S\ref{ch:secagg}. The cryptographic correctness invariant is the
only contribution we claim from the SecAgg layer; that survives the
simplifications.

\paragraph{``Why not MIMIC-IV / PhysioNet challenge data / X?''}
MIT-BIH is the canonical ECG benchmark; all other datasets we considered
(MIMIC-IV vital signs, the PhysioNet 2020 12-lead challenge) introduce
either feature heterogeneity or scale that the framework does not
currently model. The framework is dataset-agnostic at the loader level
(\texttt{src/data/}); switching to MIMIC-IV is a one-file extension and
is listed in \S\ref{sec:future}.

\section{Practical implications for deployment}

\begin{itemize}[leftmargin=*]
  \item \textbf{Run SecAgg by default.} Free utility cost (\S\ref{sec:f5.2}),
  meaningful threat coverage (honest-but-curious server,
  per-update tracing), and the correctness check is one unit test.
  \item \textbf{DP-SGD on imbalanced clinical data is non-trivial.}
  The headline-accuracy number can mask total class-prior collapse.
  Always report majority-class baseline and macro-F1.
  \item \textbf{Compression is free in the post-DP regime.} If the model
  is already at the privacy floor, top-$k$ + int8 compression buys
  bandwidth at no further utility cost. If the model is \emph{not} at
  the floor, compression introduces additional approximation error and
  must be measured explicitly.
\end{itemize}

\section{What we deliberately did not test}

\begin{itemize}[leftmargin=*]
  \item Byzantine / poisoning attacks. The threat model is
  honest-but-curious server plus model-recipient adversary.
  \item Real network-level performance. The Flower simulation is in-process.
  \item Patient-level privacy guarantees (vs.\ record-level). Each MIT-BIH
  record corresponds to one patient in our partition, but the DP
  $\varepsilon$ is per-record, not per-patient. Patient-level $\varepsilon$
  via group privacy is an extension we list but did not implement.
\end{itemize}

\section{Future work}\label{sec:future}

\begin{enumerate}[leftmargin=*]
  \item \textbf{Multi-seed runs across all five configurations} to put error
  bars on Table~\ref{tab:detection-grid}. Roughly 4 hours of GPU time per
  3-seed sweep; limiting factor is the full-pipeline run (30 min/seed).
  \item \textbf{Class-weighted DP-SGD} to test the proposed mitigation
  in \S\ref{sec:negotiating} on the records 200--234 slice, where
  DP-SGD alone clears the cliff but the full pipeline does not.
  Records 200--234 already lifted the rare-class signal; class-weighted
  loss is the next dial.
  \item \textbf{Modular-integer SecAgg} to remove the float32 cancellation
  noise observed in \S\ref{sec:f5.2} and bring the SecAgg overhead to
  exactly zero.
  \item \textbf{Byzantine-robust aggregation} (Krum, median, trimmed mean)
  composed with the existing privacy stack.
  \item \textbf{MIMIC-IV vital-signs port} to test whether the class-prior
  collapse generalises to a different healthcare modality with different
  label distribution.
  \item \textbf{Diffie--Hellman key exchange + Shamir secret-sharing} to
  upgrade the SecAgg prototype to a dropout-tolerant deployment-grade
  protocol.
\end{enumerate}

\chapter{Conclusion}\label{ch:concl}

We built a layered-defense privacy-preserving federated learning framework for
ECG arrhythmia classification on MIT-BIH, with three independently switchable
defenses (DP-SGD, pairwise-mask SecAgg, top-$k$~+~int8 compression) on top of
FedAvg. Five configurations were benchmarked on real PhysioNet data with four
simulated wearable IoT clients.

Three findings define the thesis:

\begin{enumerate}[leftmargin=*]
  \item \textbf{Federation is free} on this workload --- vanilla FedAvg
  matches centralized accuracy within 0.02\,pp. The federated decomposition
  introduces no measurable utility cost.
  \item \textbf{Secure aggregation is free} up to float32 mask-cancellation
  noise --- 0.20\,pp gap vs.\ vanilla FedAvg, well below run-to-run
  variance. The cryptographic correctness property is encoded as a unit test
  that runs in CI; the protocol is mechanically verified, not just argued.
  \item \textbf{DP-SGD collapses to majority-class prediction} on this
  imbalanced workload. The final accuracy at $\varepsilon = 5$ is 0.9711,
  which equals the prevalence of class \emph{N} (Normal) to four decimals
  --- the model is, with statistical certainty, the majority-class
  predictor. Adding compression on top of DP does not degrade further
  because there is nothing left to degrade. A follow-up on records
  200--234 (\S\ref{sec:f6.3}) shows that the collapse is
  \emph{dataset-conditional}: on moderate imbalance (82\,\% N) DP-SGD
  alone retains a $+6.37$\,pp signal over the majority baseline, but
  the full pipeline still collapses at $\varepsilon = 8$. The
  sharpness of the cliff is a property of the workload imbalance; its
  existence is a property of the privacy stack composition.
\end{enumerate}

The single recommendation this thesis makes to the
privacy-preserving-FL-for-healthcare community: report majority-class
baseline and macro-F1 alongside overall accuracy. The privacy/utility
tradeoff curve in any imbalanced-data regime is meaningless without them;
in our experiments it is the only thing standing between an honest
``DP collapsed the model'' and a reportable ``DP costs 1.7\,pp.''

The protocol and infrastructure to do so are contributed in \texttt{src/} and
\texttt{experiments/}, all driven by committed YAML configurations. Every
number in this document can be reproduced with \texttt{python -m
experiments.run\_*} on the corresponding config.

Future work (\S\ref{sec:future}) is multi-seed validation, records 200--234,
modular-integer SecAgg, and a MIMIC-IV port.

\begin{thebibliography}{99}
\bibitem{mcmahan2017fedavg}
H. Brendan McMahan, Eider Moore, Daniel Ramage, Seth Hampson, and Blaise
Ag{\"u}era y Arcas.
\emph{Communication-Efficient Learning of Deep Networks from Decentralized
Data}.
AISTATS 2017, pp.\ 1273--1282.

\bibitem{abadi2016dpsgd}
Mart{\'i}n Abadi, Andy Chu, Ian Goodfellow, H.~Brendan McMahan, Ilya Mironov,
Kunal Talwar, and Li Zhang.
\emph{Deep Learning with Differential Privacy}.
ACM CCS 2016, pp.\ 308--318.

\bibitem{mironov2017rdp}
Ilya Mironov.
\emph{R{\'e}nyi Differential Privacy}.
IEEE Computer Security Foundations Symposium (CSF) 2017, pp.\ 263--275.

\bibitem{bonawitz2017secagg}
Keith Bonawitz, Vladimir Ivanov, Ben Kreuter, Antonio Marcedone, H.~Brendan
McMahan, Sarvar Patel, Daniel Ramage, Aaron Segal, and Karn Seth.
\emph{Practical Secure Aggregation for Privacy-Preserving Machine Learning}.
ACM CCS 2017, pp.\ 1175--1191.

\bibitem{yousefpour2021opacus}
Ashkan Yousefpour, Igor Shilov, Alexandre Sablayrolles, Davide Testuggine,
Karthik Prasad, Mani Malek, John Nguyen, Sayan Ghosh, Akash Bharadwaj, Jessica
Zhao, Graham Cormode, and Ilya Mironov.
\emph{Opacus: User-Friendly Differential Privacy Library in PyTorch}.
NeurIPS Privacy in ML Workshop, 2021.

\bibitem{beutel2020flower}
Daniel J.\ Beutel, Taner Topal, Akhil Mathur, Xinchi Qiu, Javier Fernandez-Marques,
Yan Gao, Lorenzo Sani, Kwing Hei Li, Titouan Parcollet, Pedro Porto Buarque
de Gusm{\~a}o, and Nicholas D.\ Lane.
\emph{Flower: A Friendly Federated Learning Research Framework}.
arXiv:2007.14390, 2020.

\bibitem{moody2001mitbih}
George B.\ Moody and Roger G.\ Mark.
\emph{The Impact of the MIT-BIH Arrhythmia Database}.
IEEE Engineering in Medicine and Biology Magazine, 20(3):45--50, 2001.

\bibitem{aami1998ec57}
Association for the Advancement of Medical Instrumentation.
\emph{Testing and Reporting Performance Results of Cardiac Rhythm and
ST Segment Measurement Algorithms (ANSI/AAMI EC57)}, 1998.

\bibitem{kachuee2018ecg}
Mohammad Kachuee, Shayan Fazeli, and Majid Sarrafzadeh.
\emph{ECG Heartbeat Classification: A Deep Transferable Representation}.
IEEE International Conference on Healthcare Informatics (ICHI) 2018,
pp.\ 443--444.

\bibitem{li2020fedprox}
Tian Li, Anit Kumar Sahu, Manzil Zaheer, Maziar Sanjabi, Ameet Talwalkar,
and Virginia Smith.
\emph{Federated Optimization in Heterogeneous Networks}.
MLSys 2020.

\bibitem{konecny2016fedlearning}
Jakub Kone{\v{c}}n{\'y}, H.~Brendan McMahan, Felix X.\ Yu, Peter Richt{\'a}rik,
Ananda Theertha Suresh, and Dave Bacon.
\emph{Federated Learning: Strategies for Improving Communication Efficiency}.
NeurIPS Workshop on Private Multi-Party Machine Learning, 2016.

\bibitem{kim2022rigorous}
Siwon Kim, Kukjin Choi, Hyun-Soo Choi, Byunghan Lee, and Sungroh Yoon.
\emph{Towards a Rigorous Evaluation of Time-series Anomaly Detection}.
AAAI 2022, pp.\ 7194--7201.

\bibitem{shokri2017membership}
Reza Shokri, Marco Stronati, Congzheng Song, and Vitaly Shmatikov.
\emph{Membership Inference Attacks Against Machine Learning Models}.
IEEE Symposium on Security and Privacy (S\&P) 2017, pp.\ 3--18.

\bibitem{kairouz2021advances}
Peter Kairouz, H.~Brendan McMahan, Brendan Avent, Aur{\'e}lien Bellet,
Mehdi Bennis, et~al.
\emph{Advances and Open Problems in Federated Learning}.
Foundations and Trends in Machine Learning, 14(1--2):1--210, 2021.

\bibitem{dwork2006dp}
Cynthia Dwork.
\emph{Differential Privacy}.
ICALP 2006, pp.\ 1--12.

\bibitem{gopi2021prv}
Sivakanth Gopi, Yin Tat Lee, and Lukas Wutschitz.
\emph{Numerical Composition of Differential Privacy}.
NeurIPS 2021.

\bibitem{mcmahan2018dpfedavg}
H.~Brendan McMahan, Daniel Ramage, Kunal Talwar, and Li Zhang.
\emph{Learning Differentially Private Recurrent Language Models}.
ICLR 2018.

\bibitem{wei2020dpfl}
Kang Wei, Jun Li, Ming Ding, Chuan Ma, Howard~H. Yang, Farhad Farokhi,
Shi Jin, Tony Q.~S. Quek, and H.~Vincent Poor.
\emph{Federated Learning with Differential Privacy: Algorithms and
Performance Analysis}.
IEEE Transactions on Information Forensics and Security, 15:3454--3469,
2020.

\bibitem{bell2020secaggplus}
James Henry Bell, Kallista A.\ Bonawitz, Adri{\`a} Gasc{\'o}n, Tancr{\`e}de
Lepoint, and Mariana Raykova.
\emph{Secure Single-Server Aggregation with (Poly)Logarithmic Overhead}.
ACM CCS 2020, pp.\ 1253--1269.

\bibitem{so2022lightsecagg}
Jinhyun So, Chaoyang He, Chien-Sheng Yang, Songze Li, Qian Yu,
Ramy E.\ Ali, Basak Guler, and Salman Avestimehr.
\emph{LightSecAgg: A Lightweight and Versatile Design for Secure
Aggregation in Federated Learning}.
MLSys 2022.

\bibitem{bernstein2018signsgd}
Jeremy Bernstein, Yu-Xiang Wang, Kamyar Azizzadenesheli, and
Anima Anandkumar.
\emph{signSGD: Compressed Optimisation for Non-Convex Problems}.
ICML 2018, pp.\ 560--569.

\bibitem{aji2017topk}
Alham Fikri Aji and Kenneth Heafield.
\emph{Sparse Communication for Distributed Gradient Descent}.
EMNLP 2017, pp.\ 440--445.

\bibitem{rieke2020future}
Nicola Rieke, Jonny Hancox, Wenqi Li, Fausto Milletari, Holger~R. Roth,
Shadi Albarqouni, Spyridon Bakas, Mathieu~N. Galtier, Bennett~A. Landman,
Klaus Maier-Hein, et~al.
\emph{The Future of Digital Health with Federated Learning}.
npj Digital Medicine, 3:119, 2020.

\bibitem{yeom2018membership}
Samuel Yeom, Irene Giacomelli, Matt Fredrikson, and Somesh Jha.
\emph{Privacy Risk in Machine Learning: Analyzing the Connection to
Overfitting}.
IEEE Computer Security Foundations Symposium (CSF) 2018, pp.\ 268--282.

\bibitem{geiping2020inverting}
Jonas Geiping, Hartmut Bauermeister, Hannah Dr{\"o}ge, and
Michael Moeller.
\emph{Inverting Gradients --- How Easy Is It to Break Privacy in
Federated Learning?}
NeurIPS 2020.

\bibitem{melis2019exploit}
Luca Melis, Congzheng Song, Emiliano De~Cristofaro, and Vitaly Shmatikov.
\emph{Exploiting Unintended Feature Leakage in Collaborative Learning}.
IEEE Symposium on Security and Privacy (S\&P) 2019, pp.\ 691--706.

\bibitem{blanchard2017krum}
Peva Blanchard, El~Mahdi El~Mhamdi, Rachid Guerraoui, and Julien Stainer.
\emph{Machine Learning with Adversaries: Byzantine Tolerant Gradient
Descent}.
NeurIPS 2017.

\bibitem{yin2018byzantine}
Dong Yin, Yudong Chen, Kannan Ramchandran, and Peter Bartlett.
\emph{Byzantine-Robust Distributed Learning: Towards Optimal Statistical
Rates}.
ICML 2018, pp.\ 5650--5659.

\end{thebibliography}

\end{document}
